\documentclass[12pt]{article}
\usepackage{amsmath, amssymb, geometry, enumitem}
\geometry{margin=1in}
\setlength{\parskip}{0.6em}
\setlength{\parindent}{0pt}

\begin{document}

\begin{center}
\Large \textbf{ECON 300 -- Labour Economics} \\[0.3em]
\large Step-by-Step Solutions: Quiz on Chapters 11 and 16 \\[0.3em]
\large Immigration and Unemployment
\end{center}

\section*{Question 1: Immigration Composition and Selection Policy}

Total immigrants:
\[
420{,}000
\]

Categories:
\begin{itemize}
    \item Economic immigrants: \(189{,}000\)
    \item Family-class immigrants: \(147{,}000\)
    \item Refugees and humanitarian immigrants: \(84{,}000\)
\end{itemize}

\subsection*{(a) Percentage share of each category}

The formula is:

\[
\text{Percentage share}=\frac{\text{category size}}{\text{total immigration}}\times 100
\]

\textbf{Economic immigrants:}
\[
\frac{189{,}000}{420{,}000}\times 100 = 45\%
\]

\[
\boxed{45\%}
\]

\textbf{Family-class immigrants:}
\[
\frac{147{,}000}{420{,}000}\times 100 = 35\%
\]

\[
\boxed{35\%}
\]

\textbf{Refugees and humanitarian immigrants:}
\[
\frac{84{,}000}{420{,}000}\times 100 = 20\%
\]

\[
\boxed{20\%}
\]

\subsection*{(b) New category sizes and new total}

\textbf{Economic immigrants increase by 18\%:}
\[
189{,}000(1.18)=223{,}020
\]

\[
\boxed{223{,}020}
\]

\textbf{Family-class immigrants decrease by 8\%:}
\[
147{,}000(1-0.08)=147{,}000(0.92)=135{,}240
\]

\[
\boxed{135{,}240}
\]

\textbf{Refugees and humanitarian immigrants unchanged:}
\[
84{,}000
\]

\[
\boxed{84{,}000}
\]

\textbf{New total number of immigrants:}
\[
223{,}020 + 135{,}240 + 84{,}000 = 442{,}260
\]

\[
\boxed{442{,}260}
\]

\subsection*{(c) New percentage share of economic immigrants}

\[
\frac{223{,}020}{442{,}260}\times 100 \approx 50.42\%
\]

\[
\boxed{50.42\%}
\]

\vspace{1em}

\section*{Question 2: Unemployment, Participation, and Employment Rates}

Working-age population:
\[
6{,}800{,}000
\]

Labour force participation rate:
\[
67\%
\]

Unemployment rate:
\[
6.5\%
\]

\subsection*{(a) Size of the labour force}

\[
\text{Labour force} = 0.67 \times 6{,}800{,}000
\]

\[
= 4{,}556{,}000
\]

\[
\boxed{4{,}556{,}000}
\]

\subsection*{(b) Number of unemployed workers}

\[
\text{Unemployed} = 0.065 \times 4{,}556{,}000
\]

\[
= 296{,}140
\]

\[
\boxed{296{,}140}
\]

\subsection*{(c) Number of employed workers}

\[
\text{Employed} = \text{Labour force} - \text{Unemployed}
\]

\[
= 4{,}556{,}000 - 296{,}140 = 4{,}259{,}860
\]

\[
\boxed{4{,}259{,}860}
\]

\subsection*{(d) Employment rate}

\[
\text{Employment rate}=\frac{\text{Employed}}{\text{Working-age population}}\times 100
\]

\[
=\frac{4{,}259{,}860}{6{,}800{,}000}\times 100 \approx 62.65\%
\]

\[
\boxed{62.65\%}
\]

Now suppose during a recession:

\begin{itemize}
    \item \(52{,}000\) employed workers lose jobs and begin searching
    \item \(36{,}000\) unemployed workers become discouraged and leave the labour force
    \item \(18{,}000\) people previously out of the labour force enter and immediately begin searching
\end{itemize}

\subsection*{(e) New number of unemployed workers}

Start with initial unemployed:

\[
296{,}140
\]

Add workers who lose jobs and begin searching:

\[
296{,}140 + 52{,}000 = 348{,}140
\]

Subtract unemployed workers who become discouraged:

\[
348{,}140 - 36{,}000 = 312{,}140
\]

Add new entrants who begin searching:

\[
312{,}140 + 18{,}000 = 330{,}140
\]

\[
\boxed{330{,}140}
\]

\subsection*{(f) New labour force}

Start with initial labour force:

\[
4{,}556{,}000
\]

Workers who become discouraged leave the labour force:

\[
4{,}556{,}000 - 36{,}000 = 4{,}520{,}000
\]

People entering and searching join the labour force:

\[
4{,}520{,}000 + 18{,}000 = 4{,}538{,}000
\]

Note: workers who lose jobs but continue searching remain in the labour force, so they do not change the total labour force.

\[
\boxed{4{,}538{,}000}
\]

\subsection*{(g) New unemployment rate}

\[
\text{New unemployment rate}=\frac{330{,}140}{4{,}538{,}000}\times 100
\]

\[
\approx 7.27\%
\]

\[
\boxed{7.27\%}
\]

\textbf{Interpretation:}

The unemployment rate rises from \(6.5\%\) to \(7.27\%\), showing a deterioration in labour market conditions. However, it still does not fully reflect hardship because some discouraged workers left the labour force and are no longer counted as unemployed.

\vspace{1em}

\section*{Question 3: Long-Term Unemployment and Labour Market Interpretation}

Each country has a labour force of:

\[
12{,}500{,}000
\]

Unemployment rate in both countries:

\[
5.6\%
\]

\subsection*{(a) Number of unemployed workers in each country}

\[
\text{Unemployed} = 0.056 \times 12{,}500{,}000
\]

\[
= 700{,}000
\]

So each country has:

\[
\boxed{700{,}000}
\]

\subsection*{(b) Number of long-term unemployed workers}

\textbf{Country A:}

\[
0.14 \times 700{,}000 = 98{,}000
\]

\[
\boxed{98{,}000}
\]

\textbf{Country B:}

\[
0.41 \times 700{,}000 = 287{,}000
\]

\[
\boxed{287{,}000}
\]

\subsection*{(c) Long-term unemployment rate as a percentage of the labour force}

Formula:

\[
\text{Long-term unemployment rate}=\frac{\text{Long-term unemployed}}{\text{Labour force}}\times 100
\]

\textbf{Country A:}

\[
\frac{98{,}000}{12{,}500{,}000}\times 100 = 0.784\%
\]

\[
\boxed{0.78\%}
\]

\textbf{Country B:}

\[
\frac{287{,}000}{12{,}500{,}000}\times 100 = 2.296\%
\]

\[
\boxed{2.30\%}
\]

\subsection*{(d) Why Country B may still have a more serious labour market problem}

Although both countries have the same overall unemployment rate of \(5.6\%\), Country B has a much larger share of long-term unemployed workers. This suggests more serious labour market difficulties because long-term unemployment often leads to:

\begin{itemize}
    \item skill deterioration,
    \item weaker re-employment prospects,
    \item greater financial hardship,
    \item more persistent structural problems.
\end{itemize}

So Country B likely faces a more severe unemployment problem even though the headline unemployment rate is the same.

\[
\boxed{\text{Country B has the more serious labour market problem because unemployment is more persistent.}}
\]

\vspace{1em}

\section*{Final Answers Summary}

\subsection*{Question 1}
\begin{align*}
(a)\quad & 45\%,\; 35\%,\; 20\% \\
(b)\quad & 223{,}020,\; 135{,}240,\; 84{,}000,\; 442{,}260 \\
(c)\quad & 50.42\%
\end{align*}

\subsection*{Question 2}
\begin{align*}
(a)\quad & 4{,}556{,}000 \\
(b)\quad & 296{,}140 \\
(c)\quad & 4{,}259{,}860 \\
(d)\quad & 62.65\% \\
(e)\quad & 330{,}140 \\
(f)\quad & 4{,}538{,}000 \\
(g)\quad & 7.27\%
\end{align*}

\subsection*{Question 3}
\begin{align*}
(a)\quad & 700{,}000 \text{ in each country} \\
(b)\quad & 98{,}000 \text{ in A},\; 287{,}000 \text{ in B} \\
(c)\quad & 0.78\% \text{ in A},\; 2.30\% \text{ in B}
\end{align*}

\end{document}